% Part: first-order-logic
% Chapter: tableaux
% Section: rules-and-proofs

\documentclass[../../../include/open-logic-section]{subfiles}

\begin{document}

\iftag{FOL}
      {\olfileid[pt]{fol}{tab}{rul}}
      {\olfileid[pt]{pl}{tab}{rul}}

\olsection{Regras e \usetoken{P}{tableau}}

!!^a{tableau} é um exame sistemático das maneiras possíveis de
!!a{sentence} ser verdadeira ou falsa em !!a{structure}. Os elementos
constitutivos de um tableau são !!{signed formula}s: !!{sentence}s
acrescidas de um valor de verdade usado como ``sinal'', $\True$ ou
então~$\False$. Essas !!{formula}s com sinal são dispostas em uma árvore
(que cresce para baixo).

\begin{defn}
  Uma \emph{!!{signed formula}} é um par constituído por um valor de
  verdade e !!a{sentence}, O.s.\@, uma das duas formas:
  \[
  \sFmla{\True}{!A} \text{ ou } \sFmla{\False}{!A}.
  \]
\end{defn}

Intuitivamente, podemos ler $\sFmla{\True}{!A}$ como ``$!A$ poderia ser
verdadeira'' e $\sFmla{\False}{!A}$ como ``$!A$ poderia ser falsa'' (em
alguma !!{structure}).

Cada !!{signed formula} na árvore ou é uma \emph{hipótese} (as hipóteses
são listadas bem no alto da árvore), ou é obtida a partir de
!!a{signed formula} acima dela por uma de várias regras de inferência.
Há duas regras para cada possível !!{main operator} da !!{formula}
anterior: uma para o caso em que o sinal é~$\True$ e outra para o caso em
que o sinal é~$\False$. Algumas regras permitem que a árvore se
ramifique, enquanto outras apenas acrescentam !!{signed formula}s ao
ramo. Uma regra pode ser (e muitas vezes precisa ser) aplicada não à
!!{signed formula} imediatamente anterior, mas a qualquer
!!{signed formula} no ramo que vai da raiz até o local em que a regra é
aplicada.

Um ramo é \emph{fechado} quando contém tanto $\sFmla{\True}{!A}$ quanto
$\sFmla{\False}{!A}$. Um !!{tableau} fechado é aquele em que todo ramo
está fechado. Segundo a interpretação intuitiva, qualquer ramo descreve
uma possibilidade conjunta, mas $\sFmla{\True}{!A}$ e
$\sFmla{\False}{!A}$ não são conjuntamente possíveis. Em outras palavras,
se um ramo está fechado, a possibilidade que ele descreve foi excluída.
Em particular, isso significa que um !!{tableau} fechado exclui todas as
possibilidades de tornar simultaneamente verdadeira toda hipótese da
forma $\sFmla{\True}{!A}$ e falsa toda hipótese da
forma~$\sFmla{\False}{!A}$.

Um !!{tableau} fechado \emph{para $!A$} é um !!{tableau} fechado com
raiz~$\sFmla{\False}{!A}$. Se existe um !!{tableau} fechado desse tipo,
todas as possibilidades de~$!A$ ser falsa foram excluídas; O.s.\@, $!A$
deve ser verdadeira em toda !!{structure}.

\end{document}
